% =====================================================================
%  DRIFTGUARD — Appendix: the multi-LLM oracle benchmark, in full.
%  Narrative appendix (paper voice), self-contained: benchmark design,
%  headline accuracy + trust calibration, the two paradoxes, reasoning
%  styles, CoT case studies, the anomaly catalogue + dataset noise floor,
%  the Llama true-positive halt, and the matched-budget cross-check.
%  Sourced from reboot/docs/Oracle_Benchmark.tex + Appendix.tex (facts).
%  Requires: booktabs, multirow, enumitem. Defines \label{app:bench}.
%  References fig:llama-halt (defined in safety.tex) and the safety-layer
%  labels ssec:verifier, ssec:verifier-gate, ssec:auditor, sec:noise.
% =====================================================================

\appendix

\section{The Multi-LLM Oracle Benchmark}
\label{app:bench}

The main paper fixes one oracle and shows the loop recovers with it. Two
deployment questions remain, and they are the kind a reviewer is right to press:
\emph{which} LLM should serve as the oracle, and does the safety machinery of
Section~\ref{sec:safety} --- the verifier and the health monitor --- generalise
beyond the single model it was calibrated on? We answer both by running the
\emph{identical} canonical loop across a panel of providers and then turning the
verifier back on their outputs as a label-free auditor. The short version is
counter-intuitive on every axis: the choice of model matters far less than the
structure of the prompt, the verifier's own confidence signals must be read
against the grain, and the providers that fail do so in exactly the ways the
health monitor is built to catch. This appendix gives the full evidence.

\subsection{Benchmark design}
\label{app:bench-design}

Every cell runs the same canonical configuration --- drift-anchored selection,
prompt~B, periodic trigger, verifier off in-loop (audited retrospectively),
health monitor on --- over the full PhreshPhish stream at a budget-friendly
$K{=}20$ ($\approx 25$ rounds, $\approx 500$ calls per provider,
$\approx$\$15--30 for the panel). Only the provider varies, so any difference is
attributable to the model and not to the pipeline around it. Eight providers were
attempted; two returned too few usable labels to audit and became the
broken-oracle case studies of Appendix~\ref{app:bench-llama}, leaving six audited.
Because $K{=}20$ sits below the $K{=}50$ of the main results, the benchmark is
designed to measure provider \emph{behaviour} --- label accuracy, trust
calibration, hallucination, reasoning style --- rather than absolute recovery, so
the accuracy we report is oracle accuracy on the AL-queried boundary samples, the
cleanest quantity that is comparable across models.

\subsection{Accuracy and trust calibration}
\label{app:bench-acc}

Table~\ref{tab:bench-acc} reports accuracy and trust calibration over the full
$25$-round budget, and the first thing to read in it is a trap. Llama-3.3-70B tops
the raw accuracy column, but not because it is the best oracle: it \emph{halts}
at round~18 (Appendix~\ref{app:bench-llama}) and therefore never faces the hard
late rounds --- the Oct~22--24 phishing cluster --- that drag down every other
model's average. On a rounds-matched comparison, recomputing each provider over
the same first $18$ rounds Llama completed, Llama falls to third ($0.87$) while
Mistral-Large-3 ($0.90$) and DeepSeek-V4-Flash ($0.87$) lead, and every full-$25$
model gains $3$--$6$ points once those late rounds are excluded.

One reading survives both framings, and it is the result that matters: the top of
the ranking is \emph{not} the frontier closed-source tier. Open-weights and
European models lead at matched budget, the GPT-5 family sits mid-pack, and the
gaps are small --- $8.9$ points best-to-worst on the raw column, and far less once
budgets are matched. The decomposition says why. Phishing accuracy is uniformly
high across all six providers ($0.95$--$1.00$); the \emph{entire} cross-model
variance lives on the benign class ($0.73$--$0.85$) --- the same class asymmetry
the verifier of Section~\ref{ssec:verifier} is built around. The practical
implication is that the structured prompt, not the model, governs oracle quality,
which is exactly what makes the loop portable across providers. The matched-budget
end-to-end cross-check of Appendix~\ref{app:bench-smallbench} corroborates it from
the other direction, with providers clustering within $\approx 4$ points.

Trust calibration is the column to watch for safety. The Pearson $\rho$ between
verifier trust and correctness ranges from $0.43$ (GPT-5.1) down to $0.16$
(Mistral): a high-accuracy oracle whose trust is \emph{decoupled} from correctness
is the dangerous deployment case, because the verifier then reports confidence
everywhere and flags none of the errors. This is the first hint that trust is an
observability signal, not a gate --- a point the next two paradoxes settle.

\begin{table}[h]
\centering
\caption{Oracle accuracy on AL-queried samples and trust calibration, over the
full $25$-round budget; $n$ is the number of usable (queried-and-parsed) labels.
Phishing accuracy is uniformly high; cross-model variance is confined to the
benign class. $^{\dagger}$Llama halted at round~18
(Appendix~\ref{app:bench-llama}), so its $n$ --- and hence its full-run accuracy
--- covers an earlier, easier subset and is \emph{not} budget-comparable; on a
rounds-matched basis it is mid-pack (see text).}
\label{tab:bench-acc}
\small
\begin{tabular}{@{}lrcccc@{}}
\toprule
Provider & $n$ & Acc & Acc$_{\text{ph}}$ & Acc$_{\text{bn}}$ & $\rho_{\tau,c}$ \\
\midrule
Llama-3.3-70B$^{\dagger}$ & 323 & 0.864 & 1.000 & 0.854 & 0.38 \\
Mistral-Large-3   & 490 & 0.859 & 0.972 & 0.827 & 0.16 \\
GPT-5.4-mini      & 481 & 0.823 & 0.951 & 0.780 & 0.31 \\
GPT-5.4           & 482 & 0.813 & 0.982 & 0.764 & 0.37 \\
DeepSeek-V4-Flash & 479 & 0.808 & 0.979 & 0.764 & 0.30 \\
GPT-5.1           & 484 & 0.775 & 1.000 & 0.730 & \textbf{0.43} \\
\bottomrule
\end{tabular}
\end{table}

\subsection{Two paradoxes: why verifier trust is observability, not a gate}
\label{app:bench-paradox}

The cross-model audit exposes two robust paradoxes that, taken together, replicate
the gating negative result of Section~\ref{ssec:verifier-gate} on an entirely
different (unstructured) benchmark --- and explain why DRIFTGUARD uses the
verifier as an auditor (Section~\ref{ssec:auditor}) rather than a per-label
filter.

\paragraph{The hallucination paradox.} Samples on which a provider hallucinates
--- the verifier refutes at least one claim --- are \emph{more} accurate than
clean ones, for all six providers (Table~\ref{tab:bench-halluc}; Fisher exact
odds ratios $1.8$--$4.7$, $p<0.001$ for five of six, Llama the only borderline
case at $p=0.10$). The mechanism is not fabrication but verbosity: phishing pages
afford many concrete indicators, so a model enumerates $\sim$8--10 of them ---
making at least one refutable almost certain --- yet labels correctly, whereas
benign-looking phishing (wallet drainers, parked redirects) affords few claims and
is precisely where it mislabels. The contradicted-claim count therefore co-measures
how phish-rich a sample is, not how truthful the model is, and is unsafe as a
per-label quality gate.

\begin{table}[h]
\centering
\caption{Hallucination paradox: accuracy is higher on hallucinated samples than on
clean ones for all six providers. Llama is the only non-significant case
($p=0.10$).}
\label{tab:bench-halluc}
\small
\begin{tabular}{@{}lcccc@{}}
\toprule
Provider & acc$_{\text{hall}}$ & acc$_{\text{clean}}$ & $\Delta$ & OR \\
\midrule
GPT-5.1           & 0.898 & 0.650 & $+24.8$ & 4.72 \\
GPT-5.4           & 0.908 & 0.678 & $+23.0$ & 4.69 \\
DeepSeek-V4-Flash & 0.870 & 0.667 & $+20.4$ & 3.36 \\
Mistral-Large-3   & 0.931 & 0.742 & $+18.9$ & 4.69 \\
GPT-5.4-mini      & 0.900 & 0.725 & $+17.5$ & 3.41 \\
Llama-3.3-70B     & 0.899 & 0.834 & $+6.4$  & 1.76 \\
\bottomrule
\end{tabular}
\end{table}

\paragraph{The very-high-trust paradox.} Bucketing labels into trust quartiles,
accuracy climbs from Q1 to Q3 but then \emph{falls} in the top quartile Q4 --- for
every provider (Table~\ref{tab:bench-trust}; e.g.\ DeepSeek Q3 $0.95\!\to$ Q4
$0.69$). A near-saturated trust score requires many indicators, all verified or
unverifiable, which is easiest to achieve on samples with a clear \emph{surface}
match to the predicates --- a set that happens to include many benign pages with
login forms, hidden iframes, or redirects. The verifier audits claim-grounding,
not sample-level ambiguity, so a rule ``accept if $\tau\geq Q4$'' would
preferentially admit a worse-than-average slice. Both paradoxes point the same
way: trust belongs in the health monitor as a cross-round aggregate, never on an
individual label.

\begin{table}[h]
\centering
\caption{Trust-quartile accuracy. Q4 (highest trust) is below Q3 for every
provider.}
\label{tab:bench-trust}
\small
\begin{tabular}{@{}lcccc@{}}
\toprule
Provider & Q1 & Q2 & Q3 & Q4 \\
\midrule
DeepSeek-V4-Flash & 0.585 & 0.949 & 0.952 & 0.692 \\
GPT-5.4           & 0.554 & 0.942 & 0.975 & 0.785 \\
GPT-5.4-mini      & 0.587 & 0.950 & 0.959 & 0.798 \\
GPT-5.1           & 0.438 & 0.926 & 0.960 & 0.771 \\
Mistral-Large-3   & 0.766 & 0.940 & 0.968 & 0.765 \\
Llama-3.3-70B     & 0.654 & 0.976 & 0.916 & 0.895 \\
\bottomrule
\end{tabular}
\end{table}

\subsection{Reasoning styles across families}
\label{app:bench-style}

Providers differ in reasoning \emph{style} far more than in quality
(Table~\ref{tab:bench-style}), and the difference is legible in the verifier's
own resolution counts. Llama is \emph{concrete}: short, generic,
predicate-matchable indicators, which earns it the highest verify rate ($0.51$)
but the least forensic value. The GPT-5.4 family is \emph{abstract}: verbose
reasoning ($\sim$9.5 indicators per call) that cites HTML and OSINT fields by name
and is consequently classed \emph{unverifiable} most often ($0.58$). Two
consequences follow that bear directly on prompt design. First, only GPT-5.4 and
Llama actually exploit the prompt's OSINT dossier as named facts; DeepSeek
consumes it silently and Mistral cites ``domain age'' without ever quoting the
number. Second, GPT-5.4-mini pairs the highest self-confidence ($0.96$) with the
highest hallucination rate ($0.12$): its most decisive calls are also its most
fabricated --- something the verifier's trust captures (it lands mid-tier) but the
model's own confidence does not. The upshot is that a lower hallucination rate or
a higher verify rate is partly just a \emph{less specific} claim style, so neither
is a truthfulness ranking --- only a verbosity signature.

\begin{table}[h]
\centering
\caption{Reasoning-style profile. The verify/unverifiable split is a verbosity
signature, not a quality ranking; confidence and hallucination are not
self-correcting.}
\label{tab:bench-style}
\small
\begin{tabular}{@{}lcccc@{}}
\toprule
Provider & conf & verify & unver & contr \\
\midrule
Llama-3.3-70B     & 0.92 & \textbf{0.513} & 0.370 & 0.117 \\
GPT-5.4-mini      & 0.96 & 0.436 & 0.439 & \textbf{0.124} \\
GPT-5.1           & 0.91 & 0.395 & 0.514 & 0.091 \\
DeepSeek-V4-Flash & 0.96 & 0.357 & 0.545 & 0.098 \\
Mistral-Large-3   & 0.97 & 0.356 & 0.550 & 0.094 \\
GPT-5.4           & 0.93 & 0.339 & \textbf{0.582} & 0.079 \\
\bottomrule
\end{tabular}
\end{table}

\subsection{Chain-of-thought case studies}
\label{app:bench-cot}

Inspecting cached reasoning on shared sample IDs makes the styles --- and the
failure modes --- concrete.

\begin{itemize}[leftmargin=1.2em,itemsep=2pt,topsep=2pt]
\item \textbf{Textbook agreement} (sample 933,
\texttt{trezorrloggin.blogspot.com}). Ground truth phishing; all six providers
correct. GPT-5.4 cites the literal lookalike string and the iframe origin; Llama
simply lists ``typosquatting domain, hidden iframe'' --- the easy branch of the
AL pool, where any reasonable oracle suffices.
\item \textbf{Unanimous against ground truth} (sample 1430,
\texttt{imdb.com/list/\dots}). Labelled phishing in PP, yet all six call it benign
($0.93$--$0.99$ confidence), correctly reading the IMDb branding, the Amazon
script chain, and the 30-year domain age. This is a probable \emph{dataset}
mislabel --- the canonical \texttt{cross\_model\_split} pattern of
Appendix~\ref{app:bench-anom}, not a model error.
\item \textbf{A true cross-model blind spot} (sample 1679,
\texttt{gestaodecadastrosatualizar.blogspot.com}). Ground truth phishing; four of
five anchor on the Blogspot $404$ payload (``no login form'') and call benign,
missing that the URL reads ``update account records'' in Portuguese. Only Mistral
catches it, via its ``free Blogspot subdomain commonly abused'' prior --- a
reminder that the oracle helps only when at least one provider's prior is right.
\item \textbf{Lucky hallucination} (sample 869, a Roblox/Robux scam). Mistral
labels phishing (correct) but four of its supporting claims land as
\emph{contradicted}; trust drops, the label stays right. For retraining the label
is all that matters --- the fabricated support is irrelevant, which is precisely
why filtering on hallucination would have discarded a correct label.
\end{itemize}

\subsection{Anomaly catalogue and the dataset noise floor}
\label{app:bench-anom}

A retrospective extractor sorted $994$ audited verdicts into six operational
error classes (Table~\ref{tab:bench-anom}), and the two \emph{empty} classes are
as informative as the populated ones. \texttt{groundless\_phish}$=0$ across
$\sim$2{,}800 rows means every phishing verdict carries a named, checkable
indicator --- so the ``the LLM just post-hoc rationalises'' critique simply does
not apply to this oracle --- and \texttt{format\_failure}$=0$ confirms the JSON
contract holds for the healthy providers. The most consequential populated class
is \texttt{cross\_model\_split} ($243$): unanimous-LLM verdicts that contradict
the ground truth, on manual triage an estimated $70$--$80\%$ \emph{dataset}
mislabels --- legitimate IMDb lists, the YouTube homepage, German classifieds, all
labelled phishing in PP yet unanimously benign with strong OSINT legitimacy. This
is a genuine benchmark label-noise floor: it means part of the gap between the
LLM-AL result ($0.80$) and the ground-truth ceiling ($0.87$) is the dataset's own
noise, not the oracle's, and the true oracle ceiling is therefore higher than the
nominal numbers suggest.

\begin{table}[h]
\centering
\caption{Anomaly catalogue ($994$ verdicts). Two empty classes validate the
prompt; \texttt{cross\_model\_split} exposes a dataset label-noise floor.}
\label{tab:bench-anom}
\small
\begin{tabular}{@{}lrl@{}}
\toprule
Class & count & meaning \\
\midrule
inverse\_conservatism & 392 & benign call, evidence says phish \\
confident\_wrong      & 310 & high-confidence wrong verdict \\
cross\_model\_split   & 243 & unanimous LLMs vs.\ GT \\
lucky\_halluc         & 49  & correct label, refuted support \\
groundless\_phish     & \textbf{0} & phish with no named indicator \\
format\_failure       & \textbf{0} & unparseable output (healthy) \\
\bottomrule
\end{tabular}
\end{table}

\subsection{Health-monitor validation on real broken oracles}
\label{app:bench-llama}

The benchmark also stress-tests the health monitor of
Section~\ref{ssec:safety-layer} on the three oracles that failed for real, and
each isolates a different point. \textbf{Grok-4} returned \emph{zero} usable
labels in its first two rounds ($0/20$), stalled at the round-0 baseline, and is
excluded from the audit. \textbf{Kimi-K2.6} is the case that justifies checking
health at the \emph{parse} layer rather than the network layer: in a prior full
run all $1{,}232$ of its API calls returned HTTP~200 with zero transport errors,
yet $839$ ($68\%$) carried a body that did not parse to the schema, so only $33\%$
of queried labels were usable and the run recovered nothing ($+0.05$ points). A
network-level liveness check would have called the oracle perfectly healthy; the
monitor's parse-rate signal sees $0.33 \ll 0.80$ and halts within the first
rounds.

\textbf{Llama-3.3-70B} is the true-positive halt, and the cleanest demonstration
that the monitor is \emph{reactive}, not predictive (Figure~\ref{fig:llama-halt}).
Its parse rate is chronically marginal --- dipping repeatedly into the warn zone
and once critical ($0.75$ at round~5) --- and it finally strings together three
consecutive sub-$0.95$ rounds ($0.85$ and $0.80$ at rounds 16--17, and a third
below $0.95$ at round~18), tripping the sustained-parse rule and halting at
round~18. The monitor does not anticipate the collapse; it fires only \emph{after}
the malformed-JSON rate has stayed degraded for three rounds running. That the
downstream PP accuracy still looks healthy at the halt is both irrelevant to the
decision and invisible to it --- the monitor is label-free and watches the output
stream, never $F_{1,w}$. The halt nonetheless strictly improves the outcome
(Table~\ref{tab:llama}): running all $25$ rounds with the monitor disabled ends at
$\mathrm{PP_{test}}\,F_{1,w}=0.6115$, \emph{below} the round-0 baseline of
$0.6234$, whereas the monitored run stops at $0.6265$. Having flagged three broken
oracles while never tripping on the four healthy providers across $25$ rounds
each, the monitor is calibrated tightly enough to act and loosely enough to avoid
false alarms.

\begin{table}[h]
\centering
\caption{The monitored halt on Llama-3.3-70B is a true positive: disabling the
monitor and running to completion ends below the round-0 baseline ($0.6234$).}
\label{tab:llama}
\small
\begin{tabular}{@{}lcc@{}}
\toprule
Llama-3.3-70B run & rounds & $\mathrm{PP_{test}}\,F_{1,w}$ \\
\midrule
monitor on (halts) & 18 & \textbf{0.6265} \\
monitor off (full) & 25 & 0.6115 \\
\bottomrule
\end{tabular}
\end{table}

\subsection{Matched-budget provider cross-check}
\label{app:bench-smallbench}

Finally, to confirm that \emph{absolute} recovery --- not just per-label accuracy
--- is provider-insensitive, the loop is run end-to-end on a reduced PP subset
($|\mathrm{PP_{stream}}|{=}500$, $|\mathrm{PP_{test}}|{=}250$) at matched budget
across four providers and two selectors (Table~\ref{tab:smallbench}). The absolute
numbers are deliberately not comparable to the full-PP main results --- the
smaller stream and label budget lower the ceiling, and $\mathrm{PP_{test}}\,n{=}250$
carries $\mathrm{SE}\approx\pm2.5$~pp on $F_{1,w}$ --- so the table supports only a
\emph{relative} reading. On that reading, providers cluster within
$\approx 3$--$4$ points per selector, consistent with the benign-class-only
variance of Appendix~\ref{app:bench-acc}: under a fixed pipeline, the oracle is
interchangeable to within a few points, and the structured prompt is doing the
work.

\begin{table}[h]
\centering
\caption{Matched-budget cross-check ($\mathrm{PP_{stream}}{=}500$, seed $2025$).
Providers cluster within $\approx 4$ points per selector; absolute values are not
comparable to the full-PP main results.}
\label{tab:smallbench}
\small
\begin{tabular}{@{}llcc@{}}
\toprule
Selector & Provider & PP $F_{1,w}$ & PP $R_{\text{phish}}$ \\
\midrule
\multirow{4}{*}{drift-anchored}
 & DeepSeek-V4-Flash & 0.675 & 0.402 \\
 & Mistral-Large-3   & 0.671 & 0.418 \\
 & GPT-5.4           & 0.635 & 0.352 \\
 & GPT-5.4-mini      & 0.635 & 0.352 \\
\midrule
\multirow{4}{*}{pure-random}
 & GPT-5.4-mini      & 0.734 & 0.516 \\
 & Mistral-Large-3   & 0.717 & 0.500 \\
 & DeepSeek-V4-Flash & 0.701 & 0.467 \\
 & GPT-5.4           & 0.701 & 0.467 \\
\bottomrule
\end{tabular}
\end{table}

Taken together, the benchmark establishes the two facts the main paper needs from
it: the self-repair loop is robust to the \emph{choice} of oracle --- accuracy
spans a modest $8.9$ points across six LLMs, all of the variance on the benign
class, led by open-weights and European models --- and it is safe under oracle
\emph{failure}, with the verifier exposing two paradoxes that confirm its trust is
observability rather than a gate, and the health monitor validated on three real
broken oracles including a true-positive halt that strictly improves the outcome.
